%% HAND-DRAWN for this deck, from the ISCAS'27 introduction.
%%
%% Chen et al. 2018, four markers measured individually on the same 458
%% patients, then fused.  Individually none of them is a usable test; the
%% combination is.  This is the slide that answers "why four channels":
%% the clinically useful decision is a learned combination of markers, so
%% the natural unit of integration is four concurrent channels plus the
%% processor that combines them, not one better potentiostat.

\begin{tikzpicture}
  \begin{axis}[
    width=9.6cm, height=4.8cm,
    ybar, bar width=0.92cm,
    ymin=0, ymax=105,
    ylabel={Sensitivity [\%]},
    symbolic x coords={HE4, CA-125, CEA, CA19-9, fused},
    % explicit: xtick=data would take the FIRST plot's coordinates only
    xtick={HE4, CA-125, CEA, CA19-9, fused},
    xticklabels={HE4, CA-125, CEA, CA19-9, all},
    xticklabel style={font=\small, align=center},
    ylabel style={font=\small},
    yticklabel style={font=\small},
    ytick={0,25,50,75,100},
    grid=major, ymajorgrids=true, xmajorgrids=false,
    major grid style={line width=0.2pt, draw=black!14},
    axis line style={line width=0.4pt, draw=black!45},
    tick style={line width=0.4pt, draw=black!45},
    enlarge x limits=0.13,
    nodes near coords, nodes near coords style={font=\small\bfseries},
    every node near coord/.append style={/pgf/number format/fixed,
                                         /pgf/number format/precision=1},
  ]
    % bar shift=0pt: two \addplots under ybar are otherwise offset side by
    % side, which puts every bar off its tick.
    \addplot[draw=black!45, fill=black!22, bar shift=0pt] coordinates
      {(HE4,63.8) (CA-125,62.8) (CEA,38.8) (CA19-9,35.7)};
    \addplot[draw=mtxA!60!black, fill=mtxA!85, bar shift=0pt] coordinates {(fused,90.8)};
    % annotations in axis-relative coordinates, so they follow the axis box
    \draw[black!45, densely dashed, line width=0.6pt]
      (rel axis cs:0.805,0) -- (rel axis cs:0.805,1);
    \node[font=\scriptsize, black!55, anchor=north, align=center]
      at (rel axis cs:0.40,0.99) {measured individually,\\ same 458 patients};
  \end{axis}
\end{tikzpicture}
